\newcommand{\reperePierreFonctions}[6]{%
  \begin{center}
  \psset{xunit=#5,yunit=#6,algebraic=true,arrowsize=2.5pt 2,dotstyle=*,dotsize=3.5pt 0}
  \begin{pspicture*}(#1,#2)(#3,#4)
    \psgrid[subgriddiv=0,gridlabels=0pt,gridcolor=lightgray](0,0)(#1,#2)(#3,#4)
    \psaxes[labelFontSize=\scriptstyle,Dx=1,Dy=1,ticksize=-2pt 0]{->}(0,0)(#1,#2)(#3,#4)
}
\newcommand{\finReperePierreFonctions}{%
  \end{pspicture*}
  \end{center}
}

\exo
On considère une fonction $f$ représentée par la courbe ci-contre.

\begin{minipage}[t]{0.53\linewidth}
\begin{enumerate}
\item Donner l'ensemble de définition de $f$.
\item Lire les images de $-3$, $0$, $2$ et $5$.
\item Déterminer les antécédents de $1$ par $f$.
\item Résoudre graphiquement l'équation $f(x)=0$.
\item Encadrer le plus précisément possible les valeurs de $f(x)$ lorsque $x$ appartient à $[-3;5]$.
\end{enumerate}
\end{minipage}
\hfill
\begin{minipage}[t]{0.43\linewidth}
\reperePierreFonctions{-4}{-3}{6}{4}{0.48cm}{0.48cm}
  \pscurve[linecolor=blue,linewidth=1.7pt](-3,2)(-2,0)(0,-2)(2,1)(3,3)(5,-1)
  \psdots[linecolor=blue](-3,2)(5,-1)
  \rput(4.35,2.25){\textcolor{blue}{$C_f$}}
\finReperePierreFonctions
\end{minipage}

\exo
La courbe ci-contre est la représentation graphique d'une fonction $f$ définie sur $[-5;5]$.

\begin{minipage}[t]{0.53\linewidth}
\begin{enumerate}
\item Compléter le tableau de valeurs suivant.
\[
\begin{array}{|c|c|c|c|c|c|c|}
\hline
x&-5&-3&-1&1&3&5\\
\hline
f(x)&&&&&&\\
\hline
\end{array}
\]
\item Le point $A(2;1)$ appartient-il à $C_f$ ? Justifier.
\item Le point $B(-4;0)$ appartient-il à $C_f$ ? Justifier.
\item Traduire l'information graphique $(-1;2)\in C_f$ par une égalité faisant intervenir $f$.
\end{enumerate}
\end{minipage}
\hfill
\begin{minipage}[t]{0.43\linewidth}
\reperePierreFonctions{-6}{-2}{6}{4}{0.42cm}{0.55cm}
  \psline[linecolor=red!75!black,linewidth=1.7pt](-5,-1)(-3,0)(-1,2)(1,3)(3,1)(5,2)
  \psdots[linecolor=red!75!black](-5,-1)(-3,0)(-1,2)(1,3)(3,1)(5,2)
  \rput(4.35,3.15){\textcolor{red!75!black}{$C_f$}}
\finReperePierreFonctions
\end{minipage}

\exo
On a représenté ci-contre une fonction $f$ définie sur $[-4;6]$.

\begin{minipage}[t]{0.52\linewidth}
\begin{enumerate}
\item Résoudre graphiquement $f(x)=2$.
\item Résoudre graphiquement $f(x)=-1$.
\item Résoudre graphiquement $f(x)>0$.
\item Résoudre graphiquement $f(x)\leq 2$.
\item En déduire le tableau de signes de $f$ sur son ensemble de définition.
\end{enumerate}
\end{minipage}
\hfill
\begin{minipage}[t]{0.44\linewidth}
\reperePierreFonctions{-5}{-3}{7}{4}{0.42cm}{0.48cm}
  \pscurve[linecolor=green!45!black,linewidth=1.7pt](-4,-2)(-3,0)(-1,3)(1,1)(3,-1)(4,0)(6,2)
  \psdots[linecolor=green!45!black](-4,-2)(6,2)
  \rput(5.2,2.8){\textcolor{green!45!black}{$C_f$}}
\finReperePierreFonctions
\end{minipage}

\exo
La courbe ci-contre représente une fonction $f$ définie sur $[-2;8]$.

\begin{minipage}[t]{0.52\linewidth}
\begin{enumerate}
\item Donner les valeurs de $x$ pour lesquelles $f$ atteint un maximum local ou un minimum local.
\item Dresser le tableau de variations de $f$.
\item Comparer, sans calcul, $f(1)$ et $f(3)$.
\item Comparer, sans calcul, $f(6)$ et $f(7)$.
\item Peut-on comparer $f(0)$ et $f(6)$ en utilisant seulement le sens de variation ? Justifier.
\end{enumerate}
\end{minipage}
\hfill
\begin{minipage}[t]{0.44\linewidth}
\reperePierreFonctions{-3}{-2}{9}{5}{0.42cm}{0.48cm}
  \psline[linecolor=purple,linewidth=1.7pt](-2,1)(0,4)(3,-1)(5,2)(8,0)
  \psdots[linecolor=purple](-2,1)(0,4)(3,-1)(5,2)(8,0)
  \rput(6.9,2.55){\textcolor{purple}{$C_f$}}
\finReperePierreFonctions
\end{minipage}

\exo
Un élève a commencé à tracer la courbe représentative d'une fonction $f$ définie sur $[-4;5]$.

\begin{minipage}[t]{0.52\linewidth}
\begin{enumerate}
\item Donner l'image de chaque point marqué sur la courbe.
\item Déterminer les antécédents de $2$ par $f$.
\item Ajouter sur le graphique un point $M$ de la courbe tel que $f(1)=3$.
\item Ajouter sur le graphique un point $N$ de la courbe tel que $N(4;-1)$.
\item Expliquer pourquoi une courbe de fonction ne peut pas contenir deux points distincts ayant la même abscisse.
\end{enumerate}
\end{minipage}
\hfill
\begin{minipage}[t]{0.44\linewidth}
\reperePierreFonctions{-5}{-3}{6}{4}{0.43cm}{0.48cm}
  \psline[linecolor=teal,linewidth=1.7pt](-4,1)(-2,2)(0,0)(2,2)(5,1)
  \psdots[linecolor=teal](-4,1)(-2,2)(0,0)(2,2)(5,1)
  \rput(4,2.45){\textcolor{teal}{$C_f$}}
\finReperePierreFonctions
\end{minipage}

\exo
On considère la fonction $f$ représentée ci-dessous.

\begin{center}
\reperePierreFonctions{-7}{-4}{7}{5}{0.55cm}{0.45cm}
  \pscurve[linecolor=blue!65!black,linewidth=1.8pt](-6,3)(-5,1)(-3,-2)(-1,-3)(1,0)(2,2)(4,3)(6,-1)
  \psdots[linecolor=blue!65!black](-6,3)(6,-1)
  \rput(5.2,2.45){\textcolor{blue!65!black}{$C_f$}}
\finReperePierreFonctions
\end{center}

\begin{enumerate}
\item Donner l'ensemble de définition de $f$.
\item Donner un encadrement de $f(x)$ sur cet ensemble.
\item Déterminer le maximum de $f$ et préciser pour quelles valeurs de $x$ il est atteint.
\item Déterminer le minimum de $f$ et préciser pour quelle valeur de $x$ il est atteint.
\item Résoudre graphiquement $f(x)\geq 2$.
\end{enumerate}

\exo
Tracer, dans chaque cas, une courbe possible. Les réponses ne sont pas uniques.

\begin{enumerate}
\item La fonction $f$ est définie sur $[-3;4]$, elle est croissante sur $[-3;1]$, décroissante sur $[1;4]$, et son maximum vaut $3$.
\item La fonction $g$ est définie sur $[0;8]$, vérifie $g(0)=2$, $g(3)=-1$, $g(8)=1$, et le nombre $0$ possède exactement deux antécédents.
\end{enumerate}

\begin{center}
\begin{minipage}{0.47\linewidth}
\reperePierreFonctions{-4}{-3}{5}{4}{0.48cm}{0.48cm}
\finReperePierreFonctions
\end{minipage}
\hfill
\begin{minipage}{0.47\linewidth}
\reperePierreFonctions{-1}{-3}{9}{4}{0.43cm}{0.48cm}
\finReperePierreFonctions
\end{minipage}
\end{center}

\exo
Les courbes ci-dessous représentent deux fonctions $f$ et $g$ définies sur $[-4;5]$.

\begin{minipage}[t]{0.52\linewidth}
\begin{enumerate}
\item Résoudre graphiquement $f(x)=g(x)$.
\item Résoudre graphiquement $f(x)>g(x)$.
\item Résoudre graphiquement $f(x)\leq g(x)$.
\item Donner un intervalle sur lequel $f$ est croissante.
\item Donner un intervalle sur lequel $g$ est décroissante.
\end{enumerate}
\end{minipage}
\hfill
\begin{minipage}[t]{0.44\linewidth}
\reperePierreFonctions{-5}{-3}{6}{5}{0.43cm}{0.45cm}
  \psline[linecolor=red,linewidth=1.7pt](-4,-1)(-2,2)(0,3)(2,1)(5,4)
  \psline[linecolor=blue,linestyle=dashed,linewidth=1.7pt](-4,3)(-2,2)(0,0)(2,1)(5,-2)
  \psdots[linecolor=red](-4,-1)(-2,2)(0,3)(2,1)(5,4)
  \psdots[linecolor=blue](-4,3)(-2,2)(0,0)(2,1)(5,-2)
  \rput(4.05,4.35){\textcolor{red}{$C_f$}}
  \rput(4.05,-1.35){\textcolor{blue}{$C_g$}}
\finReperePierreFonctions
\end{minipage}

\exo
La fonction $f$ est définie sur $[-2;7]$ par la courbe ci-dessous.

\begin{minipage}[t]{0.52\linewidth}
\begin{enumerate}
\item Dresser le tableau de signes de $f$.
\item Résoudre graphiquement $f(x)<0$.
\item Résoudre graphiquement $f(x)\geq 1$.
\item Comparer $f(-1)$, $f(2)$ et $f(6)$.
\item Rédiger une phrase qui décrit les variations de $f$ sur $[-2;7]$.
\end{enumerate}
\end{minipage}
\hfill
\begin{minipage}[t]{0.44\linewidth}
\reperePierreFonctions{-3}{-3}{8}{4}{0.43cm}{0.48cm}
  \pscurve[linecolor=orange!85!black,linewidth=1.7pt](-2,-1)(-1,0)(1,2)(3,0)(5,-2)(7,3)
  \psdots[linecolor=orange!85!black](-2,-1)(7,3)
  \rput(6.1,2.3){\textcolor{orange!85!black}{$C_f$}}
\finReperePierreFonctions
\end{minipage}
